% 宇宙学（综述）
% license CCBYSA3
% type Wiki

本文根据 CC-BY-SA 协议转载翻译自维基百科\href{https://en.wikipedia.org/wiki/Cosmology}{相关文章}。

宇宙学（Cosmology，源自古希腊语 κόσμος（cosmos，意为“宇宙”或“世界”）与 λογία（logia，意为“研究”））是一门研究宇宙——即宇宙整体本质的物理学与形而上学的分支。英语单词 “cosmology” 首次出现在 1656 年托马斯·布朗特（Thomas Blount）的《Glossographia》中，其含义为“关于世界的论述”。\(^\text{[2]}\)1731 年，德国哲学家克里斯蒂安·沃尔夫（Christian Wolff）在拉丁语中使用 “cosmologia” 一词，用以指代研究物质世界一般性质的形而上学分支。\(^\text{[3]}\)宗教或神话宇宙学（religious or mythological cosmology）是一套基于神话、宗教与秘传文献及传统的信仰体系，通常与创世神话（creation myths）及末世论（eschatology）相关。在天文学的科学范畴中，宇宙学主要关注对宇宙年代学（chronology of the universe）的研究。

物理宇宙学（physical cosmology）研究可观测宇宙的起源、其大尺度结构与动力学特征，以及宇宙的最终命运，同时探讨支配这些领域的科学定律。\(^\text{[4]}\)这一学科由科学家（如天文学家与物理学家）以及哲学家（如形而上学家、物理哲学家与时空哲学家）共同研究。由于其研究范围与哲学存在交叉，物理宇宙学中的理论既可能包含科学命题，也可能包含非科学命题，且可能依赖于无法通过实验检验的假设。物理宇宙学是天文学的一个子分支，专注于研究整个宇宙。现代物理宇宙学由“大爆炸理论”（Big Bang Theory）主导，该理论试图结合观测天文学与粒子物理学。\(^\text{[5][6]}\)更具体而言，标准的宇宙参数化模型包括暗物质（dark matter）与暗能量（dark energy），统称为“ΛCDM 模型”（Lambda-CDM model）。

理论天体物理学家戴维·N·斯珀格尔（David N. Spergel）将宇宙学称为一门“历史科学”（historical science），因为“当我们仰望宇宙时，实际上是在回望过去”，这源于光速的有限性。\(^\text{[7]}\)
\subsection{学科领域}
\begin{figure}[ht]
\centering
\includegraphics[width=8cm]{./figures/3be309e51e6413de.png}
\caption{} \label{fig_CoGy_1}
\end{figure}
物理学与天体物理学通过科学的观测与实验，在塑造人类对宇宙的理解过程中发挥了核心作用。物理宇宙学（physical cosmology）通过数学与观测相结合的方式，在对整个宇宙的分析中逐步形成。普遍认为，宇宙起源于“大爆炸”（Big Bang），随后几乎瞬间经历了“宇宙暴涨”（cosmic inflation）——即空间的迅速膨胀。据推测，宇宙大约在 $13.799 \pm 0.021$ 十亿年前由此诞生。\(^\text{[8]}\)宇宙起源学（cosmogony）研究宇宙的起源，而宇宙描绘学（cosmography）则负责绘制宇宙的结构与特征。

在狄德罗（Diderot）的《百科全书》（Encyclopédie）中，宇宙学被划分为四个部分：天学（uranology，研究天体之科学）、气学（aerology，研究空气之科学）、地学（geology，研究大陆之科学）以及水学（hydrology，研究水体之科学）。\(^\text{[9]}\)

形而上宇宙学（metaphysical cosmology）亦被描述为探讨人类在宇宙中与其他存在物之间关系的学科。罗马皇帝兼哲学家马可·奥勒留（Marcus Aurelius）对此作出了典型的论述，他指出人的位置与认知密切相关：“不知世界为何物者，不知己所处之地；不知世界为何而存者，不知己为何人，亦不知世界为何物。”\(^\text{[10]}\)
\subsection{发现}
\subsubsection{物理宇宙学}
物理宇宙学是物理学与天体物理学的一个分支，研究宇宙的物理起源与演化，同时也探讨宇宙在大尺度上的性质。在其最早期的形式中，它即如今所谓的“天体力学”（celestial mechanics）——对天体运动与“天界”的研究。希腊哲学家萨摩斯的阿里斯塔克（Aristarchus of Samos）、亚里士多德（Aristotle）与托勒密（Ptolemy）提出了不同的宇宙理论。其中，托勒密的地心体系（geocentric Ptolemaic system）长期占据主导地位，直到十六世纪，尼古拉·哥白尼（Nicolaus Copernicus）以及随后约翰内斯·开普勒（Johannes Kepler）与伽利略·伽利莱（Galileo Galilei）提出日心体系（heliocentric system）。这是物理宇宙学中最著名的认识论断裂（epistemological rupture）之一。

艾萨克·牛顿（Isaac Newton）于 1687 年发表的《自然哲学的数学原理》（Philosophiæ Naturalis Principia Mathematica）首次系统描述了万有引力定律（law of universal gravitation）。该理论为开普勒行星运动定律提供了物理机制，并能够解释行星之间引力相互作用所导致的系统异常。牛顿宇宙学与其前辈理论的一个根本区别在于哥白尼原理（Copernican principle）：地球上的物体遵循与天体相同的物理定律。这一思想构成了物理宇宙学的一个重要哲学飞跃。

现代科学宇宙学通常被认为始于 1917 年，当时阿尔伯特·爱因斯坦（Albert Einstein）发表了其对广义相对论（general relativity）的最终修正论文《广义相对论的宇宙学考量》（Cosmological Considerations of the General Theory of Relativity）。\(^\text{[11]}\)（尽管这篇论文直到第一次世界大战结束后才在德国以外被广泛传播。）广义相对论的提出促使威廉·德西特（Willem de Sitter）、卡尔·史瓦西（Karl Schwarzschild）与阿瑟·爱丁顿（Arthur Eddington）等宇宙起源学家（cosmogonists）探讨其在天文学中的影响，从而极大地提升了天文学家研究遥远天体的能力。物理学家们也由此开始放弃“宇宙是静止且不变的”这一传统假设。1922 年，亚历山大·弗里德曼（Alexander Friedmann）首次提出“宇宙膨胀”概念，即宇宙中包含运动的物质并在不断扩张。

